\documentclass[../../main.tex]{subfiles}
\begin{document}

{\bf [解]}

\paragraph{(1)}

% TODO: 概要図
% TODO: 概要図
\begin{figure}[htb]
  \centering
  \begin{tikzpicture}[scale=2.2, >=stealth, every node/.style={font=\small}]

    % ============================================================
    % 1. 定数・厳密パラメータ設定
    % 半径 t = 2.0
    % 中心角 theta = 54° (L = t * theta_rad)
    % ============================================================
    \def\tVal{2.0}
    \def\thDeg{54}

    \coordinate (O) at (0, 0);
    \coordinate (P) at (\tVal, 0);
    \coordinate (Q) at ({\tVal*cos(\thDeg)}, {\tVal*sin(\thDeg)});

    % ============================================================
    % 2. 扇形 OPQ の薄い塗りつぶし
    % ============================================================
    \fill[blue!10, opacity=0.6] (O) -- (P) arc (0:\thDeg:\tVal) -- cycle;

    % ============================================================
    % 3. 座標軸
    % ============================================================
    \draw[->, thick, gray!80!black] (-0.3, 0) -- (\tVal + 0.6, 0) node[right] {$x$};
    \draw[->, thick, gray!80!black] (0, -0.3) -- (0, \tVal + 0.3) node[above] {$y$};
    \node[below left=2pt] at (O) {$O$};

    % ============================================================
    % 4. 円弧・動径・弧長 L
    % ============================================================
    % 動径 OP (x 軸上)
    \draw[line width=1.2pt, gray!30!black] (O) -- (P);

    % 動径 OQ
    \draw[line width=1.3pt, gray!30!black] (O) -- (Q)
    node[midway, above left=1pt, font=\footnotesize] {半径 $t$};

    % 弧 PQ (太い赤線で弧長 L を強調)
    \draw[line width=1.8pt, red!85!black] (P) arc (0:\thDeg:\tVal);
    \node[red!90!black, font=\normalsize\bfseries, right=2pt]
    at ({\tVal*cos(\thDeg/2)}, {\tVal*sin(\thDeg/2)}) {弧長 $L$};

    % 半径 t の補助円弧（薄い点線）
    \draw[dashed, gray!45, line width=0.8pt] (P) arc (0:85:\tVal);

    % ============================================================
    % 5. 中心角 θ = L/t
    % ============================================================
    \draw[->, thick, teal!85!black] (0.55, 0) arc (0:\thDeg:0.55);
    \node[teal!90!black, font=\footnotesize\bfseries] at (\thDeg/2:0.75) {$\theta = \frac{L}{t}$};

    % ============================================================
    % 6. 点 P, Q のプロットと座標の投影
    % ============================================================
    \fill[black] (P) circle (1.3pt) node[below=2pt] {$P(t, 0)$};
    \fill[blue!85!black] (Q) circle (1.4pt) node[above right=1pt, font=\normalsize\bfseries] {$Q$};

    % Q からの投影破線
    \draw[dotted, gray!60, line width=0.9pt] (Q) -- ({\tVal*cos(\thDeg)}, 0)
    node[below=2pt, font=\scriptsize] {$u(t) = t\cos\frac{L}{t}$};
    \draw[dotted, gray!60, line width=0.9pt] (Q) -- (0, {\tVal*sin(\thDeg)})
    node[left=2pt, font=\scriptsize] {$v(t) = t\sin\frac{L}{t}$};

  \end{tikzpicture}
  \caption{半径 $t$ の円周上を進む点 $Q$ と弧長 $L$・中心角 $\dfrac{L}{t}$ の幾何関係}
  \label{fig:circle_arc_length_Q_parametric}
\end{figure}

題意の円周は$x^2+y^2=t^2$ と半径$t$だから，この円周上を偏角$\dfrac{L}{t}$だけ進んだ点が$Q$で，
\begin{align}
  u(t) = t \cos \dfrac{L}{t}, v(t) = t \sin \dfrac{L}{t} \label{eq:1}
\end{align}
である．

\paragraph{(2)}

\cref{eq:1}の両辺を$t$で微分して
\begin{align}
  \begin{cases}
    u'(t) = \cos \dfrac{L}{t} + t \cdot \left( -\sin \dfrac{L}{t} \right) \left( -\dfrac{L}{t^2} \right) = \cos \dfrac{L}{t} + \dfrac{L}{t} \sin \dfrac{L}{t} \\
    v'(t) = \sin \dfrac{L}{t} + t \left( \cos \dfrac{L}{t} \right) \left( -\dfrac{L}{t^2} \right) = \sin \dfrac{L}{t} - \dfrac{L}{t} \cos \dfrac{L}{t}
  \end{cases}
\end{align}
だから，
\begin{align}
  \{u'(t)\}^2 + \{v'(t)\}^2
  &= 1 + \left( \dfrac{L}{t} \right)^2 + 2 \dfrac{L}{t} \left( \cos \dfrac{L}{t} \sin \dfrac{L}{t} - \sin \dfrac{L}{t} \cos \dfrac{L}{t} \right) \\
  &= 1 + \dfrac{L^2}{t^2}
\end{align}
である．$0 < a < 1$ から，
\begin{align}
  f(a) = \int_a^1 \sqrt{1 + \dfrac{L^2}{t^2}} dt = \int_a^1 \dfrac{\sqrt{t^2+L^2}}{t} dt
\end{align}
で，$x = \sqrt{t^2+L^2}$ とすると，$\dfrac{dx}{dt} = \dfrac{t}{x}$ だから $\alpha = \sqrt{a^2+L^2}$, $\beta = \sqrt{1+L^2}$ として
\begin{align}
  f(a)
  &= \int_\alpha^\beta \dfrac{x}{t} \cdot \dfrac{x}{t} dx \\
  &= \int_\alpha^\beta \dfrac{x^2}{x^2-L^2} dx \\
  &= \int_\alpha^\beta \left( 1 + \dfrac{L}{2} \left( \dfrac{1}{x-L} - \dfrac{1}{x+L} \right) \right) dx \\
  &= \left[ x + \dfrac{L}{2} \log (x-L) - \dfrac{L}{2} \log (x+L) \right]_\alpha^\beta \\
  &= (\beta - \alpha) + \dfrac{L}{2} \log \dfrac{\beta-L}{\alpha-L} - \dfrac{L}{2} \log \dfrac{\beta+L}{\alpha+L} \\
  &= \sqrt{L^2+1} - \sqrt{L^2+a^2} + \dfrac{L}{2} \log \dfrac{(\sqrt{L^2+1}-L)}{(\sqrt{L^2+a^2}-L)} \dfrac{(\sqrt{L^2+a^2}+L)}{(\sqrt{L^2+1}+L)} \label{eq:2}
\end{align}
である．

\paragraph{(3)}

$g(a)=\dfrac{f(a)}{\log a}$とおくと\cref{eq:2}より
\begin{align}
  g(a) & = \dfrac{\sqrt{L^2+1} - \sqrt{L^2+a^2}}{\log a} \\
  & \quad + \dfrac{L}{2}\left[\dfrac{\log(\beta-L)}{\log a}+\dfrac{\alpha+L}{\log a} - \dfrac{\alpha - L }{\log a} - \dfrac{\beta + L}{\log a}\right] \label{eq:3}
\end{align}
である．各項$a\to +0$の極限をとると
\begin{align}
  & \lim_{a\to +0} \dfrac{\sqrt{L^2+1} - \sqrt{L^2+a^2}}{\log a} = 0 \label{eq:4} \\
  & \lim_{a\to +0} \dfrac{\log (\beta \pm L)}{\log a} = 0 \label{eq:5} \\
  & \lim_{a\to +0} \dfrac{\log (\alpha + L)}{\log a} =  0 \label{eq:6}\\
  & \lim_{a\to +0} \dfrac{\log (\alpha - L)}{\log a} = \lim_{a\to +0} \dfrac{\log (\sqrt{L^2+a^2} - L)}{\log a} \\
  & \quad \quad \quad = \lim_{a\to +0} \dfrac{\log a^2 - \log (\sqrt{L^2+a^2} + L)}{\log a} = 2 \label{eq:7}
\end{align}
だから，\cref{eq:4,eq:5,eq:6,eq:7}を\cref{eq:3}に代入して求める極限値は
\begin{align}
  \lim_{a\to +0}  \dfrac{f(a)}{\log a} = - \dfrac{L}{2} \cdot 2 = -L
\end{align}
である．

\end{document}
